% Everything after a percent sign is a comment.
% The document class sets the overall layout.
\documentclass[a4paper,11pt]{article}

% The preamble: settings for the whole document.
\title{Counting Words}
\author{A. Geek}
\date{}  % empty: print no date

\begin{document}  % the text starts here
\maketitle        % print title and author

% An empty line starts a new paragraph.
This is a \emph{short} sample document. It shows
\textbf{bold} text, \emph{italic} text and
\verb|inline code|.

\section{Introduction}  % numbered automatically
Counting words sounds simple. It is simple,
if you know your tools.

\subsection{Why count words}
\begin{itemize}  % a bullet list
  \item Essays often have a word limit.
  \item Editors pay by the word.
  \item It is fun.
\end{itemize}

\section{How to do it}
\begin{enumerate}  % a numbered list
  \item Save your text as a plain text file.
  \item Open a terminal.
  \item Run one of the commands below.
\end{enumerate}

% {lll}: three columns, all left aligned.
% & separates cells, \\ ends a row.
\begin{tabular}{lll}
  \hline
  Tool   & Command                 & Notes    \\
  \hline
  wc     & \verb|wc -w file.txt|    & fast     \\
  Python & \verb|python3 count.py| & flexible \\
  \hline
\end{tabular}

\subsection{A small program}
% verbatim prints the text exactly as typed.
\begin{verbatim}
import sys

words = open(sys.argv[1]).read().split()
print(len(words), "words")
if len(words) < 100:
    print("a short text & a quick read")
\end{verbatim}

\end{document}
